\documentclass{book}
\usepackage{ctex}
\usepackage[cmex10]{amsmath}
\usepackage{xcolor}
\usepackage{mathtools}
\usepackage{booktabs}
\DeclareRobustCommand{\bluebold}[1]{\textbf{\textcolor[rgb]{0.31,0.63,0.86}{#1}}}
\usepackage[
    pdfpagemode=UseOutlines,
    bookmarks=true,
    bookmarksopen=true,
    colorlinks=true,
    linkcolor=black,
    citecolor=black,
    urlcolor=black
]{hyperref}
\usepackage{bookmark}

\raggedbottom

\begin{document}
\tableofcontents
\clearpage


\chapter{Fluid statics 流体静力学}

\section{Properties of fluids 流体性质}

Continuity of fluids 流体连续性

流体是由无数彼此相连、彼此无间隙的\bluebold{质点}组成，是一种\bluebold{连续性}介质

\bluebold{质点}：大量分子或原子组成的流体微团
虚拟的，保持流体宏观特性的最小流体微元

\bluebold{尺寸}：远小于管道、设备尺寸，远大于分子自由程

\section{Compressibility 可压缩性}

体积压缩系数

\begin{equation}
\beta = \frac{1}{\rho} \left( \frac{\partial \rho}{\partial p} \right)_T
\label{eq:isothermal-compressibility}
\end{equation}

\begin{itemize}
    \item 下标 T 代表等温条件，是压缩系数必不可少的标识；
\end{itemize}

or

\begin{equation}
\beta = \frac{1}{\rho} \frac{d\rho}{dp}
\end{equation}

beta 越大, 流体越容易被压缩

\subsection{体积压缩系数}

下标 T 代表等温条件，是压缩系数必不可少的标识；

\section{无定形性}

没有固定的形状，形状随容器而变

\section*{notes}

体积压缩系数

\chapter{密度和压强}

\section{密度 \texorpdfstring{$\rho$}{密度 ρ}}

单位体积流体所具有的质量，$\mathrm{kg/m^3}$ or$\quad \mathrm{kg\cdot m^{-3}}$

\begin{equation}
\rho = \frac{m}{V} = f(T,p)
\label{eq:density-def}
\end{equation}

由公式 \ref{eq:density-def} 可知密度与温压相关。

可压缩流体，$\rho = f(T,p)$

不可压缩流体，$\rho \neq f(p)$, $\rho = f(T)$

\section{压强}

垂直作用在单位流体表面积上的法向力 $p = \frac{dF}{dA}$, 

标准大气压：

$\mathrm{1\,atm} = \mathrm{101325\,Pa} = \mathrm{760\,mmHg} = \mathrm{10.33\,mH_2O}$

\bluebold{绝对压强}：真实压强，以绝对真空为基准

\bluebold{表压强和真空度}：当地大气压为基准

\begin{equation}
p_\text{表} = p_\text{绝} - p_\alpha
\end{equation}

\begin{equation}
p_\text{真} = p_\alpha - p_\text{绝}
\end{equation}

表压强和真空度互为相反数

\chapter{流体静力学基本方程式}

\section{流体静力学基本方程式}

\begin{equation}
\mathrm{p_1} + \mathrm{\rho g z_1} = \mathrm{p_2} + \mathrm{\rho g z_2}
\end{equation}

\begin{equation}
\frac{p_1}{\rho} + \mathrm{g z_1} = \frac{p_2}{\rho} + \mathrm{g z_2}
\end{equation}

\begin{equation}
\mathrm{\rho_2} = \mathrm{\rho_1} + \mathrm{\rho gh}
\end{equation}

等压面：静止的、连续的、不可压缩的、同一种流体内部，处于同一水平面上各点的压强相等。

\section{静力学基本方程式的应用}

U形管压差计，指示剂A与被测液体B不互溶，$\mathrm{\rho_A} > \mathrm{\rho_B}$

静力学应用，关键是寻找等压面

\begin{equation}
p_\text{a} = p'_\text{a}
\end{equation}

\begin{equation}
p_\text{a} = p_1 + \rho_\mathrm{B} g (m + R)
\end{equation}

\begin{equation}
p'_\text{a} = p_2 + \rho_\mathrm{B} g m + \rho_\mathrm{A} g R
\end{equation}

第一个等式，是第二、三个相等的联系平衡式

\begin{equation}
p_1 + \rho_\mathrm{B} g (m + R) = p_2 + \rho_\mathrm{B} g m + \rho_A g R
\end{equation}

最终得到

\begin{equation}
p_1 - p_2 = (\rho_\mathrm{A} - \rho_\mathrm{B}) g R
\end{equation}

\end{document}